\section{Gêmeos Digitais em Prótese Dentária}

A reconstrução protética abrange uma miríade de cenários biomecânicos e estéticos, variando de coroas unitárias a reabilitações completas. A incorporação de gêmeos digitais nesses processos introduz um caráter preditivo que transcende o simples escaneamento tridimensional e a confecção mecanizada.

\subsection{Prótese Fixa e Comportamento de Materiais Cerâmicos}

A prótese fixa --- englobando coroas, pontes e facetas --- foi uma das primeiras áreas da odontologia a adotar o fluxo de trabalho digital. Contudo, o que se limitava ao Desenho Assistido por Computador (CAD) e à Manufatura Assistida por Computador (CAM) evoluiu para um ecossistema sofisticado com a ontologia dos gêmeos digitais. O modelo virtual do dente preparado é enriquecido com informações críticas: espessura da dentina remanescente, proximidade com a câmara pulpar e as propriedades elásticas (módulo de Young e coeficiente de Poisson) do material restaurador selecionado \cite{li2025impacts}.

Simulações de Análise de Elementos Finitos (FEA) são executadas diretamente sobre esse modelo paramétrico para predizer a distribuição de tensões de von Mises e as tensões principais máximas na interface dente--restauração sob cargas mastigatórias cíclicas. Essa predição de fratura é imperativa para coroas posteriores confeccionadas em materiais frágeis, como dissilicato de lítio e zircônia monolítica. Embora ostentem notável tenacidade à fratura e resistência à flexão, a susceptibilidade a defeitos de usinagem e à propagação subcrítica de trincas (\textit{slow crack growth}) pode precipitar falhas catastróficas ao longo da vida útil em função \cite{zhang2024recursive}.

\destaque{Gêmeos digitais que internalizam modelos de fadiga amparados na mecânica da fratura linear-elástica oferecem a capacidade ímpar de estimar a sobrevida da restauração, acusando cenários de risco antes da cimentação definitiva.}

O planejamento da macrotextura oclusal baseia-se em articuladores virtuais totalmente dinâmicos. O gêmeo digital simula os movimentos mandibulares com base em registros cinemáticos extraídos de axiografias digitais ou rastreamento óptico. Interações de contato prematuro são interceptadas algoritmicamente, permitindo ajustes automáticos da morfologia da coroa para garantir uma desoclusão fisiológica \cite{roy2025applications}.

\subsection{Prótese sobre Implante e Distribuição de Tensões}

O planejamento protético sobre implantes osteointegráveis beneficia-se intrinsecamente da abordagem dos gêmeos digitais, a qual amálgama em um único modelo vetorial a coordenada espacial do implante, o perfil de emergência, a geometria do componente protético (\textit{abutment}) e a restauração final. Essa consolidação é instrumental em reabilitações extensas, em que variáveis como o paralelismo dos eixos de inserção, a conformação estética dos tecidos peri-implantares e a mecânica das conexões prótese-implante ditam o prognóstico a longo prazo \cite{shen2024effectiveness}.

A análise de tensão perimplantar por métodos computacionais figura entre as aplicações mais validadas. O modelo de elementos finitos estendido do complexo coroa--pilar--implante--osso permite investigar a transferência de deformações micrométricas ao osso cortical e trabecular sob vetores de força oblíquos e axiais. Em cenários idênticos, o clínico pode mensurar quantitativamente o impacto amortecedor de diferentes materiais (e.g., coroas em resina de polieteretercetona -- PEEK, versus metalocerâmica) \cite{menon2026innovations}.

Essas plataformas simulam também o impacto de diferentes níveis ósseos e conexões protéticas (cone morse versus hexágono externo), reduzindo substancialmente a incidência de afrouxamento de parafusos de retenção e mitigando fenômenos de saucerização (perda óssea crestal) ao otimizar o selamento biológico e mecânico.

\subsection{Prótese Total e Parcial Removível}

Frequentemente negligenciadas nas discussões de ponta em odontologia digital, as próteses totais (PT) e as próteses parciais removíveis (PPR) encontram na adoção dos gêmeos digitais um vetor de intensa renovação tecnológica. O arcabouço virtual de uma prótese total transcende a simples malha superficial da mucosa de revestimento; ele incorpora as propriedades viscoelásticas do tecido mole subjacente e a relação cêntrica neuromuscular do paciente \cite{olawade2026advancing}.

A modelagem de retenção e estabilidade durante o ciclo mastigatório envolve o cálculo de vetores de deslocamento decorrentes das forças não-axiais. \textit{Softwares} especializados modelam a pressão de fluido no biofilme salivar sob a base da prótese, computando as forças de adesão, coesão e sucção capilar. Isso propicia a otimização algorítmica das bordas periféricas da prótese, visando à obtenção do vedamento periférico perfeito sem invasão das zonas de inserção muscular, algo especialmente crítico nas bases mandibulares, onde a área de suporte é diminuta.

\subsection{Reabilitação Oral Complexa e Dimensão Vertical}

Reabilitações que reestruturam todo o arranjo oclusal do paciente -- englobando desgaste erosivo crônico, atrição severa e perda de Dimensão Vertical de Oclusão (DVO) -- representam o apogeu da aplicabilidade dos gêmeos digitais. A tessitura de todos os elementos em um "paciente virtual" permite o reestabelecimento do plano oclusal, o cálculo milimétrico da DVO ideal e a harmonização estética de longo curso \cite{aisoc2026realising}.

A modificação da DVO repercute na dinâmica das Articulações Temporomandibulares (ATMs) e nos comprimentos de repouso dos músculos masseter e temporal. O gêmeo digital atua como um laboratório de provas in vivo simulado, testando o impacto das modificações nos torques articulares. Os guias de preparo e os \textit{mock-ups} de teste são gerados não de maneira estática, mas derivados da cinemática mandibular ótima, minimizando tensões musculares aberrantes e salvaguardando a durabilidade da reabilitação \cite{li2024aligner}.
